\section{Introduction}
\label{vii:sec:introduccion}

La question gravitationnelle qui organise ce travail est de savoir comment une structure
qui conserve l'histoire de ses opérations peut déterminer un
courant matériel et, par son intermédiaire, une contribution à l'évolution de
l'espace-temps. La réponse se construit du transport des routes
jusqu'à la variation de l'action : le résultat central est une amplitude
torsionnelle explicite, accompagnée d'un secteur de positivité conservée et
d'une solution cosmologique globale dans la réalisation de poussière. La thèse
relie la mémoire arithmétique à une réponse géométrique déterminée par des
opérations ; le retour de phase conserve une périodicité observable
tandis que l'état accumule l'orientation, la retenue et la profondeur.

On introduit à cet effet un noyau formel mathématique appelé holographie modulaire triadique, ci-après HMT. \emph{All Possible Paths}, ci-après APP, est une structure arithmétique de chemins sur \((\mathbb Z/9\mathbb Z)^2\). Son adjacence est décrite au moyen du graphe de Cayley additif de générateurs cardinaux \(\pm(1,0)\) et \(\pm(0,1)\), et sa réalisation cellulaire périodique est toroïdale. L'accumulation additive, la composition multiplicative et leurs évaluations maintiennent des fonctions distinctes au sein de cette structure. TRIT est l'unité de signature orientée \(\{-1,0,+1\}\), qui type le régime local de l'opération. Le \emph{Topological Phase Kernel}, ci-après TPK, constitue sa dynamique de sélection, de transport, de mémoire et de prolongement. L'état enrichi et la structure discrète conjointe du continu conservent les feuillets et les données résiduelles à travers ces opérations. La mécanique dimensionnelle réalise ses sorties en sections d'action, en grandeurs géométriques et en observables physiques ; les valeurs internes transmettent à chaque étape les relations qui les déterminent.

La question centrale est de savoir comment la mémoire transportée intervient
dans l'évolution gravitationnelle. La dilution, le rebond et sa persistance
décrivent les conséquences cosmologiques d'une amplitude qui doit d'abord
être construite à partir de l'action matérielle. La première relation s'exprime
au moyen de deux lecteurs d'une même route,
\begin{equation}
 \frac{a_n}{a_{\mathrm{ref}}}=\varphi^{n/3},
 \qquad M_n=\varphi^{-2n},
 \qquad
 M_n=\left(\frac{a_n}{a_{\mathrm{ref}}}\right)^{-6}.
 \label{vii:eq:programa-dilucion}
\end{equation}
Ici, \(a_{\mathrm{ref}}>0\) fixe la normalisation du facteur d'échelle.
La puissance exprime la relation entre la croissance de la réalisation
spatiale et le caractère quadratique du lecteur de mémoire. La contribution
gravitationnelle requiert en outre de déterminer le courant et l'amplitude qui
accompagnent cette dépendance d'échelle. Les deux constructions occupent une
section propre avant leur introduction dans les équations cosmologiques.

Le premier mouvement de l'article réunit construction et géométrie. La
décomposition du transport distingue la mémoire symétrique et la circulation
orientée. La représentation du groupoïde des routes en géométrie de Cartan
se développe avec la composition, l'inversion, la covariance et le raffinement.
L'action, le paramètre de Barbero--Immirzi, la constante gravitationnelle et
l'unité surfacique conservent leurs antécédents structurels ; leur réutilisation
transmet les conditions qui individualisent chaque section.

En particulier, la section~\ref{g26:sec:rigidez-radial-es} construit
d'abord la longueur \(L\), compose son lecteur positif avec le même
caractère qui détermine l'énergie et obtient le coefficient gravitationnel.
L'horloge circulaire reçoit alors
\(t_0=\pi_{\rm HMT}L/(54c)\) ; son rayon \(R_{108}\) coïncide avec
\(L\) dans cette réalisation. L'ordre de la démonstration distingue la
composition longitudinale, la réponse opératoire et l'évaluation en
unités physiques. Les règles longitudinale et métrique sont des compositions
explicitées dans ce développement ; le théorème prouve la rigidité des
réalisations qui les conservent, avec l'identification radiale déclarée.

Le deuxième mouvement développe l'amplitude, l'évolution et la persistance. Le courant d'écran est porté à une différentielle par rapport à la connexion au moyen de l'incidence transportée et de l'extension minimisante. La réalisation spinorielle est développée à son tour à partir des matrices internes de Clifford et de son action matérielle. L'action d'extension minimisante et l'action spinorielle sont deux fonctionnelles distinctes : la première dépend en général non linéairement de la connexion, tandis que la seconde est affine en la contorsion lorsque la cotétrade et le champ spinoriel sont fixés. Leurs variations sont effectuées sur les arguments propres à chaque fonctionnelle ; leur composition éventuelle exige de conserver cette dépendance. La normalisation traciale fixe la lecture de densité correspondante et reste distinguée d'une somme extensive. La variation de la connexion détermine la réponse torsionnelle de l'action minimale. Lorsque la cotétrade et les champs matériels évoluent, la torsion algébriquement associée à leur courant de spin change également.

L'amplitude du secteur spinoriel est déterminée par
\begin{equation}
 s_0[\varrho,V_c]
 =\frac{3\kappa c^2\hbar^2}{16}
   \frac{\gamma^2}{1+\gamma^2}
   \frac{\operatorname{Tr}
    (\varrho\widehat{\mathscr Q}_{\rm sp})}{V_c^2},
 \qquad \kappa=\frac{8\pi G}{c^4}.
 \label{vii:eq:programa-amplitud}
\end{equation}
Ici, \(\varrho\) est l'état matériel normalisé, \(V_c\) le volume
comobile et \(\widehat{\mathscr Q}_{\rm sp}\) l'observable quartique
construit à partir du courant. Le
théorème~\ref{vii:thm:amplitud-material-explicita} dérive la formule et
le corollaire~\ref{vii:cor:dominio-positivo-conservado} présente des états
dans lesquels l'amplitude est positive et reste constante. La dépendance
du second moment conserve de l'information sur l'état même lorsque la
moyenne du courant s'annule. La donnée matérielle et son volume fixent une
réalisation ; les constantes de couplage procèdent des sections
structurelles construites au préalable.

L'élimination de la torsion est également traitée pour l'action d'
extension minimale, dont la dépendance à l'égard de la connexion est, en
général, non linéaire. Sa différentielle détermine l'équation stationnaire ;
le hessien contrôle la continuation locale et une identité intégrale
détermine l'action effective. Cette formulation permet de conserver la
dépendance complète du minimum lors des variations.

La formulation du rebond conserve les densités positives, le signe de
la contribution torsionnelle et l'inégalité des exposants de la réalisation
homogène. Sa persistance locale se démontre au moyen de deux contrôles
séparés : une borne uniforme de la perturbation conserve le changement de
signe aux extrémités d'un intervalle ; une borne uniforme de sa dérivée
conserve la monotonie et la transversalité. Les réalisations des
sélecteurs cosmologiques sont énoncées avec leurs propres domaines, ainsi que les
invariants et bilans qu'elles produisent.

La solution plate de poussière avec \(\Lambda_{\mathrm{eff}}=0\)
fournit en outre un développement global :
\[
 a(t)=a_{\rm b}\left(1+
       \frac{(t-t_{\rm b})^2}{\tau^2}\right)^{1/3},
 \qquad
 a_{\rm b}^3=\frac{s_0}{\rho_0},\qquad
 \tau^2=\frac{s_0}{6\pi G\rho_0^2},
 \quad c=1.
\]
La preuve établit la régularité de la courbure et la complétude des
géodésiques causales de Levi--Civita dans cette réalisation homogène.
Son intégration et la persistance locale
du seuil décrivent deux propriétés complémentaires, avec leurs
domaines respectifs explicites.

La conservation du tenseur résiduel se formule dans la coordinatisation métrique sélectionnée, avec des couplages constants sur celle-ci et une source matérielle covariamment conservée. Sa décomposition en trace et en partie sans trace permet d'identifier l'échange entre les deux secteurs ; la constance de la trace caractérise le cas de conservation séparée.

Le troisième mouvement étudie les horizons et la réponse relationnelle. L'état
réduit, les registres accessibles et la purification permettent de décrire
quelle information une lecture représente et quelle information demeure dans les
corrélations de l'état conjoint. La définition opérationnelle de
l'observateur comprend l'orientation et la trivialisation du transport,
la restriction aux observables et la projection sur le sous-système. La réponse
auto-inertielle est reliée à ces opérations par la factorisation
par la trace globale et le défaut géométrique de projection.

\Needspace{24\baselineskip}
Les bilans d'horizon conservent également les conventions d'
aire, d'énergie, de température et d'état qu'ils utilisent.
Les sections~\ref{viii:subsec:conversion-termica-calendario-horizonte}--\ref{viii:subsec:seccion-termometrica-registro}
composent le quantum aréal avec l'énergie de décision du calendrier et
obtiennent la conversion entre les températures chronologique et cosmologique.
L'état réduit conserve les occupations des deux secteurs
angulaires. À rayon fixé, sa réponse énergétique a pour hamiltonien
\(\mathsf H_R=\tau_R\Delta_0\mathsf D_\partial\) ; ses variations
produisent la première loi, tandis que les différences finies retiennent
le terme d'entropie relative. Ce développement construit une section
thermométrique relative et y établit une preuve d'égalité entre
applications linéaires de conversion. La provenance sectorielle reste
accessible tant dans l'aire que dans le coût énergétique.
La propagation avancée et retardée pour des sources à support fini se
compose sur le transport unitaire des routes. L'égalité de leurs lectures à la frontière s'
exprime au moyen d'une restriction linéaire des sources ; sa projection
détermine la réponse variationnelle dans ce domaine.

L'unité de l'argument réside dans la relation entre source, courant, amplitude et géométrie. La mémoire conserve une information que la lecture de phase n'épuise pas ; sa réalisation matérielle produit un second moment capable de contribuer à la dynamique même lorsque le courant moyen est nul. L'évolution cosmologique exprime cette contribution à l'échelle globale dans le secteur résolu, et les opérations de l'observateur précisent quelle partie de l'information reste accessible. La structure, l'équation et sa solution se trouvent ainsi reliées au sein de réalisations dont les domaines sont explicites.
